\documentclass{../../plantilla/raysac-informe}

\producto{transbank-SDK}
\titulo{Informe de Diagramas de Arquitectura y Flujos}
\titulocorto{Informe de Diagramas}
\version{1.0.0}
\fecha{2026-11-22}
\clasificacion{Uso interno}
\autor{raysac}

\begin{document}

\portada

\begin{controlversiones}
  \cvfila{1.0}{2026-11-22}{raysac}{Versión inicial con la colección completa de diagramas de arquitectura, secuencias, estados, dominio y despliegue.}
\end{controlversiones}

\avisolegalestandar

\clearpage
\preliminares

\section{Introducción y marco de modelado}

El presente documento recopila y describe exhaustivamente el conjunto de diagramas arquitectónicos,
estructurales y de comportamiento del proyecto \texttt{transbank-SDK}. Este proyecto proporciona un
SDK no oficial en el lenguaje de programación Rust para la comunicación bidireccional con terminales
de punto de venta (POS) físicos de Transbank (modelos POS Integrado y POS Autoservicio) a través de
interfaces serie/USB y puentes de red TCP, integrando adicionalmente un simulador de terminal y una
interfaz de línea de comandos (CLI).

\subsection{Notación y estándar visual}

Todos los diagramas han sido elaborados utilizando PlantUML bajo el estándar visual compartido del
ecosistema raysac (\texttt{docs/plantilla/estilo.puml}). Las fuentes oficiales residen en
\texttt{docs/diagramas/fuentes/*.puml}, siendo compiladas de forma determinista a representaciones
vectoriales SVG y mapas de bits PNG de alta resolución mediante el script \texttt{docs/build-docs.sh}.

Se emplean convenciones formales de modelado:
\begin{itemize}
  \item \textbf{C4 Model (Niveles 1 y 2-3):} para delimitar las fronteras del sistema, actores humanos,
    sistemas externos y la descomposición interna en componentes de software y bibliotecas.
  \item \textbf{UML 2.5:} para diagramas de clases de dominio, máquinas de estado de enlace, diagramas
    de actividad de procesos de negocio, casos de uso y diagramas de secuencia temporal.
  \item \textbf{Estereotipo \texttt{<<planificado>>}:} utilizado para destacar visualmente con estilo
    atenuado componentes o flujos contemplados en el diseño pero cuya implementación final está sujeta
    a validaciones posteriores con hardware físico o tareas programadas en el roadmap.
\end{itemize}

\subsection{Catálogo de diagramas incluidos}

La Tabla~\ref{tab:catalogo-diagramas} enumera los diez diagramas que conforman el cuerpo de este informe.

\begin{longtable}{@{}p{1.2cm}p{5.8cm}p{3.2cm}p{4.2cm}@{}}
\caption{Catálogo de diagramas del proyecto transbank-SDK.}\label{tab:catalogo-diagramas}\\
\toprule
\textbf{N°} & \textbf{Diagrama} & \textbf{Tipo de modelo} & \textbf{Archivo fuente} \\
\midrule
\endfirsthead
\toprule
\textbf{N°} & \textbf{Diagrama} & \textbf{Tipo de modelo} & \textbf{Archivo fuente} \\
\midrule
\endhead
\bottomrule
\endlastfoot
01 & Contexto del sistema & C4 Nivel 1 & \texttt{contexto.puml} \\
02 & Componentes del workspace & C4 Nivel 2-3 & \texttt{componentes.puml} \\
03 & Casos de uso por actor & UML Casos de Uso & \texttt{casos-de-uso.puml} \\
04 & Clases y tipos del dominio & UML Clases & \texttt{clases-dominio.puml} \\
05 & Estados del enlace & UML Máquina de Estados & \texttt{estados-enlace.puml} \\
06 & Flujo de proceso de venta & UML Actividad & \texttt{flujo-venta.puml} \\
07 & Secuencia de venta exitosa & UML Secuencia & \texttt{secuencia-venta.puml} \\
08 & Recuperación y nunca doble cobro & UML Secuencia & \texttt{secuencia-recuperacion-venta.puml} \\
09 & Despliegue de infraestructura & UML Despliegue & \texttt{despliegue.puml} \\
10 & Ecosistema retail y facturación & Integración C4 & \texttt{ecosistema.puml} \\
\end{longtable}

\clearpage
\section{Diagrama 01: Contexto del sistema (C4 Nivel 1)}

\begin{figure}[htbp]
  \centering
  \includegraphics[width=\textwidth]{contexto.png}
  \caption{Contexto del sistema transbank-SDK (C4 Nivel 1).}
  \label{fig:contexto}
\end{figure}

\subsection{Propósito}
El diagrama de contexto (Figura~\ref{fig:contexto}) ilustra la posición de \texttt{transbank-SDK}
dentro del entorno operativo de una tienda comercial o caja de autoservicio. Delimita claramente la
responsabilidad de la biblioteca respecto al software de facturación host, los operadores humanos, el
dispositivo de hardware físico (POS) y la infraestructura bancaria adquirente operada por Transbank.

\subsection{Elementos y actores}
\begin{itemize}
  \item \textbf{Cajero / Operador:} usuario humano que opera el software de punto de venta (caja atendida).
  \item \textbf{Cliente (Pagador):} titular de la tarjeta que interactúa directamente con el terminal POS
    para insertar/acercar su tarjeta e ingresar su clave secreta (PIN).
  \item \textbf{Técnico / Integrador:} desarrollador o técnico de soporte que utiliza las herramientas CLI
    para tareas de diagnóstico, prueba de puertos y certificación.
  \item \textbf{Sistema de Facturación / Host (\texttt{rework-sistema-facturacion}):} aplicación de escritorio
    (construida con Tauri y Rust) responsable de la venta y emisión de boletas tributarias.
  \item \textbf{transbank-SDK (\texttt{transbank-pos}):} biblioteca de enlace y protocolo objeto del presente
    diseño, ejecutada en el espacio de memoria del proceso host.
  \item \textbf{Terminal POS Transbank:} dispositivo de hardware (Verifone VX520, Ingenico Desk 3500) que
    captura la tarjeta y procesa la transacción bancaria.
  \item \textbf{Red Transbank:} switch adquirente centralizado que valida fondos y emite autorizaciones.
  \item \textbf{Simulador POS (\texttt{tbk-pos-sim}):} emulador de terminal que reproduce fielmente el
    protocolo sin requerir hardware físico conectado.
  \item \textbf{Servicio de Impuestos Internos (SII):} entidad tributaria receptora de las boletas
    electrónicas tras la confirmación del pago.
\end{itemize}

\subsection{Decisiones de diseño reflejadas}
\begin{itemize}
  \item \textbf{Aislamiento PCI DSS (SEC-01):} el host y el SDK nunca acceden al número de tarjeta completo
    (PAN) ni al PIN. La captura sensible se realiza íntegramente dentro del terminal físico.
  \item \textbf{Independencia de hardware mediante simulador (ADR-10):} el sistema permite el desarrollo y
    las pruebas automáticas sin depender de un POS físico conectado.
\end{itemize}

\subsection{Guía de lectura}
El flujo comienza en la esquina superior izquierda cuando el cajero inicia la operación en la caja.
El software host delega la orden en el SDK mediante métodos de alto nivel. El SDK transforma la orden
en bytes de protocolo y los envía al POS. El cliente ingresa su PIN en el POS; el POS contacta a Transbank
y devuelve el resultado final al SDK. Una vez aprobado, el host genera la boleta electrónica ante el SII.

\clearpage
\section{Diagrama 02: Componentes del workspace (C4 Nivel 2-3)}

\begin{figure}[htbp]
  \centering
  \includegraphics[width=\textwidth]{componentes.png}
  \caption{Componentes del workspace transbank-SDK (C4 Nivel 2-3).}
  \label{fig:componentes}
\end{figure}

\subsection{Propósito}
El diagrama de componentes (Figura~\ref{fig:componentes}) detalla la arquitectura modular del workspace
Cargo y la separación en capas dentro del crate principal \texttt{transbank-pos}, el simulador
\texttt{transbank-pos-sim} y la herramienta de terminal \texttt{transbank-pos-cli}.

\subsection{Estructura de paquetes y capas}
\begin{itemize}
  \item \textbf{Capa de cliente (\texttt{client}):} contiene las estructuras \texttt{PosIntegrado},
    \texttt{PosAutoservicio} y la variante asíncrona \texttt{AsyncPosIntegrado}. Proporciona métodos
    de alto nivel con validación de argumentos y tipado estricto.
  \item \textbf{Capa de enlace (\texttt{link}):} estructura \texttt{Link<T, C>} encargada del protocolo
    de tramas: transmisión, espera de acuse de recibo (ACK/NAK), reintentos acotados, control de plazos
    (timeouts), descarte de ruido y procesamiento de tramas intermedias (0900). Comparte lógica con
    el módulo \texttt{exchange}.
  \item \textbf{Capa de protocolo sans-IO (\texttt{protocol}):} lógica pura libre de efectos colaterales de
    entrada/salida o reloj. Incluye el cálculo del checksum (\texttt{lrc}), codificación (\texttt{frame}),
    máquina de estados de decodificación (\texttt{FrameDecoder}), formateo y parseo de mensajes
    (\texttt{message}) y tipos de dominio (\texttt{types}, \texttt{codes}).
  \item \textbf{Capa de transporte (\texttt{transport}):} abstracción mediante el trait \texttt{Transport},
    con implementaciones para memoria con tiempo inyectable (\texttt{MemoryTransport}), sockets TCP
    (\texttt{TcpTransport}) y puertos serie físicos RS-232/USB (\texttt{SerialTransport}).
  \item \textbf{Crate simulador (\texttt{transbank-pos-sim}):} núcleo sans-IO (\texttt{PosSimulator}),
    inyección de escenarios de falla (\texttt{Scenario}), transporte en memoria (\texttt{SimTransport})
    y servidor de red TCP (\texttt{tbk-pos-sim}).
  \item \textbf{Crate CLI (\texttt{transbank-pos-cli}):} binario \texttt{tbk-pos} que expone todas las
    operaciones del cliente a la línea de comandos mediante \texttt{clap}.
\end{itemize}

\subsection{Decisiones de diseño reflejadas}
\begin{itemize}
  \item \textbf{Arquitectura sans-IO (ADR-01):} la lógica de análisis y generación de bytes no realiza
    llamadas al sistema operativo, permitiendo una cobertura de pruebas exhaustiva mediante pruebas basadas
    en propiedades (\texttt{proptest}) y reutilización directa en el simulador.
  \item \textbf{Inyección de reloj (ADR-09):} el enlace parametriza el tiempo mediante el trait \texttt{Clock},
    lo que permite ejecutar pruebas de timeouts de cientos de segundos en microsegundos de CPU.
\end{itemize}

\subsection{Guía de lectura}
Las peticiones fluyen de arriba hacia abajo: desde el software del comercio hacia la capa de clientes,
descendiendo al enlace que coordina la serialización de tramas en \texttt{protocol} y la emisión de
bytes sobre la interfaz \texttt{Transport}. Las respuestas recibidas ascienden en sentido inverso,
validándose su integridad antes de instanciar los tipos de dominio tipados.

\clearpage
\section{Diagrama 03: Casos de uso por actor}

\begin{figure}[htbp]
  \centering
  \includegraphics[width=\textwidth]{casos-de-uso.png}
  \caption{Diagrama de casos de uso por actor.}
  \label{fig:casos-de-uso}
\end{figure}

\subsection{Propósito}
El diagrama de casos de uso (Figura~\ref{fig:casos-de-uso}) define las capacidades funcionales que el
sistema ofrece a cada uno de sus actores humanos y automatizados, organizando las operaciones
estándar de la normativa de Transbank POS Integrado y POS Autoservicio.

\subsection{Actores y responsabilidades}
\begin{itemize}
  \item \textbf{Cajero / Operador:} usuario principal de punto de venta. Ejecuta el cobro de ventas (CU-01),
    la anulación de transacciones de crédito del día (CU-04) y la recuperación de ventas dudosas (CU-03).
  \item \textbf{Cliente (Pagador):} titular de la tarjeta bancaria. Participa en la venta presentando su medio
    de pago e ingresando su número de identificación personal (PIN) en el teclado seguro del terminal.
  \item \textbf{Administrador / Supervisor:} responsable del arqueo y cuadratura. Consulta los acumulados de
    la jornada (CU-05: Totales), recupera el listado pormenorizado (CU-06: Detalle), efectúa el cierre formal
    del terminal con el switch bancario (CU-07) e inicia la carga de llaves criptográficas (CU-08).
  \item \textbf{Desarrollador / Técnico de Soporte:} utiliza la interfaz CLI \texttt{tbk-pos} para pruebas de
    enlace en terreno (CU-10) y levanta el servidor \texttt{tbk-pos-sim} para pruebas de integración (CU-11).
  \item \textbf{Software de Caja (Host):} sistema automatizado que realiza el sondeo periódico de presencia
    (CU-09: Poll) y coordina la regla de oro de recuperación ante fallas de comunicación.
\end{itemize}

\subsection{Decisiones de diseño reflejadas}
\begin{itemize}
  \item \textbf{Extensión obligatoria de recuperación (CU-03 extend CU-01):} ante cualquier interrupción o
    vencimiento del plazo de espera durante una venta, el caso de uso de cobro se extiende forzosamente
    hacia la recuperación de la última transacción, evitando cobros duplicados al consumidor.
  \item \textbf{Operación desatendida (CU-02):} el protocolo de POS Autoservicio adapta los comandos para
    entornos desatendidos sin teclado de supervisor, gestionando la impresión externa del voucher.
\end{itemize}

\subsection{Guía de lectura}
Los actores humanos se ubican en el costado izquierdo y los automatizados en el inferior. Las relaciones
de asociación directa indican inicio de acción. La relación segmentada \texttt{<<extend>>} entre el cobro
y la recuperación subraya el mecanismo de contingencia ante fallos.

\clearpage
\section{Diagrama 04: Clases y tipos del dominio}

\begin{figure}[htbp]
  \centering
  \includegraphics[width=\textwidth]{clases-dominio.png}
  \caption{Diagrama de clases y tipos de dominio de transbank-pos.}
  \label{fig:clases-dominio}
\end{figure}

\subsection{Propósito}
El diagrama de clases (Figura~\ref{fig:clases-dominio}) presenta el modelo de tipos en memoria de la
biblioteca \texttt{transbank-pos}. Ilustra cómo el sistema aprovecha el sistema de tipos estático de
Rust para impedir estados inválidos, errores de redondeo financiero o datos malformados.

\subsection{Elementos principales}
\begin{itemize}
  \item \textbf{Newtypes de dominio (\texttt{protocol::types}):}
    \begin{itemize}
      \item \texttt{Amount}: monto expresado en números enteros de pesos chilenos (\texttt{u32}). Rango
        garantizado de 50 a 999.999.999 CLP. Prohíbe el uso de tipos de punto flotante (\texttt{f32}/\texttt{f64}).
      \item \texttt{Ticket}: identificador alfanumérico de 1 a 6 caracteres ASCII, con formateo y padding
        automático de ceros a la izquierda.
      \item \texttt{OperationNumber}: número correlativo de operación de 1 a 6 dígitos para anulaciones.
      \item \texttt{CommerceCode}: código de comercio de hasta 12 dígitos, soportando código principal (``0'').
      \item \texttt{PosDate} y \texttt{PosTime}: tipos específicos para fecha (\texttt{ddMMyyyy}) y hora
        (\texttt{HHmmss}) locales informadas por el terminal, sin supuestos erróneos de zona horaria.
      \item \texttt{CardType} y \texttt{CardBrand}: enums que categorizan crédito, débito y prepago, y marcas.
    \end{itemize}
  \item \textbf{Catálogo de códigos (\texttt{codes}):}
    \begin{itemize}
      \item \texttt{ResponseCode}: enum que modela 56 códigos de respuesta oficiales de Transbank, con métodos
        auxiliares \texttt{is\_approved()}, \texttt{requires\_close()} y \texttt{requires\_load\_keys()}.
      \item \texttt{IntermediateStatus}: estados intermedios recibidos en tramas \texttt{0900}.
    \end{itemize}
  \item \textbf{Peticiones y Respuestas (\texttt{message}):}
    \begin{itemize}
      \item \texttt{SaleRequest} y \texttt{MultiCodeSaleRequest}: objetos de entrada inmutables.
      \item \texttt{SaleResponse}: estructura integral con 21 campos que preserva el payload original
        (\texttt{raw}) e implementa redacción automática de datos sensibles en \texttt{Debug}.
    \end{itemize}
\end{itemize}

\subsection{Decisiones de diseño reflejadas}
\begin{itemize}
  \item \textbf{Rechazo financiero no es Error (ADR-03):} una tarjeta rechazada produce \texttt{Ok(SaleResponse)}
    donde \texttt{is\_approved() == false}. El tipo \texttt{Result::Err} se reserva para fallos de canal o protocolo.
  \item \textbf{Tipos no exhaustivos (\texttt{\#[non\_exhaustive]}):} permite futuras expansiones de la
    especificación oficial de Transbank sin provocar rupturas en el código de los consumidores de la API.
\end{itemize}

\subsection{Guía de lectura}
Los paquetes agrupan conceptos relacionados. Los newtypes validan sus restricciones en sus constructores
\texttt{new()} retornando \texttt{Result}. Las respuestas contienen instancias opcionales de estos tipos,
reflejando con honestidad la presencia o ausencia de campos en las tramas enviadas por el POS.

\clearpage
\section{Diagrama 05: Máquina de estados del enlace}

\begin{figure}[htbp]
  \centering
  \includegraphics[width=\textwidth]{estados-enlace.png}
  \caption{Máquina de estados de la capa de enlace Link<T>.}
  \label{fig:estados-enlace}
\end{figure}

\subsection{Propósito}
El diagrama de estados (Figura~\ref{fig:estados-enlace}) modela el ciclo de vida y los estados de
transición del enlace \texttt{Link<T>}. Describe cómo se garantiza la entrega ordenada de tramas,
el manejo estricto de reintentos, el control de plazos y la inmunidad ante ruido en el puerto serie.

\subsection{Estados y transiciones}
\begin{itemize}
  \item \textbf{Desconectado:} estado inicial antes de abrir el puerto serie o conectar el socket TCP.
  \item \textbf{Inactivo (Idle):} canal operativo listo para recibir comandos. Se ejecuta un drenado
    preventivo de entrada para descartar bytes espurios residuales.
  \item \textbf{EnviandoComando:} la trama física (STX + payload + ETX + LRC) es escrita mediante
    \texttt{Transport::write\_all}. Se incrementa el contador de intentos.
  \item \textbf{EsperandoAck:} espera activa de acuse de recibo con plazo acotado (\texttt{ack\_timeout = 2 s}).
    Si expira o llega un NAK (0x15), se reintenta el envío hasta un máximo de 3 intentos (\texttt{max\_attempts}).
    Si se agotan, transiciona a \texttt{ErrorTerminal} con \texttt{Error::NoAck}.
  \item \textbf{EsperandoRespuesta:} el POS ha tomado el control. Se espera la trama de respuesta definitiva
    bajo un temporizador global (\texttt{response\_timeout = 150 s}).
    \begin{itemize}
      \item Si llega una trama \texttt{0900} (mensaje intermedio), se responde inmediatamente con ACK (0x06),
        se emite un evento al callback de la aplicación y se continúa esperando sin reiniciar el temporizador global.
      \item Los bytes fuera de trama o basura eléctrica se descartan silenciosamente sin abortar la espera.
    \end{itemize}
  \item \textbf{VerificandoLrc:} recibida la trama completa, se recalcula el checksum XOR. Si es incorrecto,
    se emite un NAK y se tolera hasta 3 fallas (\texttt{Error::BadLrc}). Si es válido, se emite ACK al POS.
  \item \textbf{RecibiendoDetalle:} estado especializado para comandos multilínea (\texttt{0261}/\texttt{0291}),
    donde las tramas se acumulan secuencialmente hasta detectar un silencio de 500 ms (\texttt{detail\_tail}).
\end{itemize}

\subsection{Decisiones de diseño reflejadas}
\begin{itemize}
  \item \textbf{Inmutabilidad del temporizador ante tramas 0900 (LINK-04):} el plazo máximo de 150 segundos
    no se reinicia por la recepción de tramas intermedias, garantizando un límite superior predecible.
  \item \textbf{ACK inmediato antes de parsear (LINK-01):} la capa de enlace reconoce la recepción física
    correcta de la trama antes de someter su contenido a análisis léxico o sintáctico.
\end{itemize}

\subsection{Guía de lectura}
Siga las flechas desde el círculo negro superior. Observe la bifurcación crítica en \texttt{EsperandoAck}:
un comando simple (como \texttt{poll} o \texttt{set\_normal\_mode}) retorna a Inactivo tras recibir ACK,
mientras que un comando con respuesta (como \texttt{sale}) pasa a \texttt{EsperandoRespuesta}.

\clearpage
\section{Diagrama 06: Flujo de proceso de venta}

\begin{figure}[htbp]
  \centering
  \includegraphics[width=\textwidth]{flujo-venta.png}
  \caption{Diagrama de actividad del proceso de venta en transbank-pos.}
  \label{fig:flujo-venta}
\end{figure}

\subsection{Propósito}
El diagrama de actividad (Figura~\ref{fig:flujo-venta}) expone la lógica algorítmica secuencial que se
desencadena cuando una aplicación consumidora invoca el método \texttt{PosIntegrado::sale(\&req)}.

\subsection{Etapas del proceso}
\begin{enumerate}
  \item \textbf{Validación de entrada:} comprobación rigurosa de límites de \texttt{Amount} (50 a 999.999.999 CLP)
    y longitud/caracteres de \texttt{Ticket}. Rechazo temprano con \texttt{Error::InvalidArgument} sin tocar el puerto.
  \item \textbf{Codificación y encuadre:} construcción del payload en formato Transbank
    (\texttt{"0200|monto|ticket|||voucher"}), cálculo de LRC mediante XOR acumulativo y encapsulamiento entre STX y ETX.
  \item \textbf{Transmisión inicial y bucle de ACK:} ciclo de envío con hasta 3 reintentos ante NAK o timeout de 2 s.
  \item \textbf{Punto de no retorno:} una vez que el POS emite el byte ACK (0x06), el comando ha ingresado formalmente
    al núcleo transaccional del hardware. A partir de este instante, el SDK prohíbe terminantemente reenviar la trama.
  \item \textbf{Bucle de recepción de respuesta:} lectura continua de bytes con descarte de ruido. Procesamiento de
    tramas intermedias \texttt{0900} (notificando al observador) hasta detectar la trama de terminación \texttt{0210}.
  \item \textbf{Validación de integridad (LRC):} comprobación de redundancia longitudinal con hasta 3 intercambios de NAK.
  \item \textbf{Parseo y resultado de negocio:} desensamblado del payload en una estructura tipada \texttt{SaleResponse}.
    Si el código es 0, se marca aprobado; en caso contrario, se consigna el motivo de rechazo entregado por el switch.
\end{enumerate}

\subsection{Decisiones de diseño reflejadas}
\begin{itemize}
  \item \textbf{Principio nunca doble cobro (ADR-04, LINK-03):} el punto de no retorno asegura que bajo ninguna
    circunstancia de falla posterior al ACK se genere una segunda orden de venta con el mismo o distinto ticket.
  \item \textbf{Parseo honesto y tolerante (PROTO-05):} los campos vacíos o ausentes se traducen a \texttt{None},
    sin inventar valores por defecto como ceros o cadenas vacías.
\end{itemize}

\subsection{Guía de lectura}
Inicie en el nodo negro superior y siga la línea de control. Observe con especial atención el nodo de decisión
de vencimiento del temporizador de 150 segundos, el cual conduce directamente a un estado de resultado
desconocido que activa el protocolo de recuperación.

\clearpage
\section{Diagrama 07: Secuencia de venta exitosa}

\begin{figure}[htbp]
  \centering
  \includegraphics[width=\textwidth]{secuencia-venta.png}
  \caption{Diagrama de secuencia de una venta aprobada en POS Integrado.}
  \label{fig:secuencia-venta}
\end{figure}

\subsection{Propósito}
El diagrama de secuencia (Figura~\ref{fig:secuencia-venta}) detalla la interacción temporal síncrona
entre los distintos objetos del sistema durante una transacción de venta estándar exitosa, visualizando
los intercambios de control a nivel de bytes en la interfaz física.

\subsection{Intercambio de mensajes y tramas}
\begin{enumerate}
  \item El host invoca \texttt{sale\_with\_status(\&req, callback)} con un monto de \$12.000 y ticket \texttt{"T00001"}.
  \item La capa cliente construye el payload y delega la ejecución en \texttt{Link::request}.
  \item \texttt{Link} formatea los bytes serie: \texttt{<STX>0200|12000|T00001|||1<ETX><LRC>} y escribe en el transporte.
  \item El POS físico responde inmediatamente con el byte \texttt{0x06} (ACK). La trama fue aceptada.
  \item El terminal solicita la interacción del cliente y envía una trama intermedia:
    \texttt{<STX>0900|01<ETX><LRC>} (``Esperando tarjeta'').
  \item \texttt{Link} emite ACK (0x06) al POS y despacha el estado al callback registrado por el host.
  \item El cliente inserta el chip e introduce su PIN en el teclado del terminal.
  \item El POS envía una segunda trama \texttt{0900|03} (``Solicitando autorización'') que se acusa con ACK.
  \item El POS transmite la solicitud financiera por su conexión de red (Ethernet/3G) al switch de Transbank.
  \item Transbank autoriza el cobro (código 00, código de autorización \texttt{"123456"}).
  \item El POS envía por serie la trama final: \texttt{<STX>0210|00|597029414300|SIM01|T00001|123456|12000|...<ETX><LRC>}.
  \item \texttt{Link} comprueba el checksum, emite el ACK final (0x06) y retorna el payload crudo al cliente.
  \item El cliente analiza el payload, genera \texttt{SaleResponse} y el POS imprime el comprobante fiscal.
\end{enumerate}

\subsection{Decisiones de diseño reflejadas}
\begin{itemize}
  \item \textbf{Visibilidad de estados intermedios sin polling (API-02):} el callback síncrono permite
    a la interfaz de usuario host reflejar en tiempo real mensajes como ``Ingrese tarjeta'' o ``Procesando pago''.
  \item \textbf{ACK bilateral riguroso (PROTO-01):} cada trama enviada o recibida es acusada formalmente con
    un byte de control ACK (0x06), cumpliendo de forma estricta la especificación de Transbank.
\end{itemize}

\subsection{Guía de lectura}
Las líneas de vida verticales representan las entidades activas. El tiempo avanza hacia abajo. Las flechas
continuas representan llamadas a métodos o transferencias de tramas completas, mientras que las flechas
segmentadas indican acuses de recibo (ACK) o retornos de estructuras de datos.

\clearpage
\section{Diagrama 08: Secuencia de recuperación (Nunca doble cobro)}

\begin{figure}[htbp]
  \centering
  \includegraphics[width=\textwidth]{secuencia-recuperacion-venta.png}
  \caption{Diagrama de secuencia del procedimiento de recuperación ante timeout.}
  \label{fig:secuencia-recuperacion-venta}
\end{figure}

\subsection{Propósito}
El diagrama de secuencia de recuperación (Figura~\ref{fig:secuencia-recuperacion-venta}) ilustra el
comportamiento del sistema ante el peor escenario de falla operativa: un corte de comunicación, desconexión
de cable o vencimiento del plazo de 150 segundos luego de que el comando de cobro fue acusado con éxito
por el POS.

\subsection{Protocolo de contingencia y reconciliación}
\begin{enumerate}
  \item El software host solicita una venta de \$15.000 asociada a la boleta interna \texttt{"BOL101"}.
  \item El comando \texttt{0200} es transmitido y el terminal POS responde con ACK (0x06). La venta está en proceso.
  \item Transcurren 150 segundos sin que el terminal entregue una respuesta definitiva (por ejemplo, congestión
    severa de red adquirente, corte del bus USB o reinicio intempestivo del terminal).
  \item \texttt{Link} agota el temporizador y retorna \texttt{Err(Error::ResponseTimeout \{ waited: 150s \})}.
  \item \textbf{Regla Mandatoria (ADR-04):} el SDK prohíbe reintentar el comando. El estado de la transacción
    bancaria es rigurosamente DESCONOCIDO (pudo haberse autorizado o rechazado internamente).
  \item El software host invoca el método \texttt{PosIntegrado::last\_sale()} (comando Transbank \texttt{0250}).
  \item El POS devuelve la última transacción asentada en su memoria no volátil mediante la trama \texttt{0260}.
  \item \textbf{Reconciliación estricta en el Host:}
    \begin{itemize}
      \item \textbf{Caso A (Coincidencia):} si el \texttt{ticket} devuelto es \texttt{"BOL101"}, el monto es
        exactamente \$15.000 y el código de respuesta es \texttt{Approved (00)}, el software de caja da por
        cobrada la venta exitosamente. Se evita un cobro duplicado al cliente.
      \item \textbf{Caso B (Discrepancia o Rechazo):} si el ticket corresponde a una venta anterior o el código
        indica rechazo financiero, se concluye con certeza que el cobro no se aplicó. El cajero es notificado
        y puede proceder con un nuevo intento bajo un nuevo número de ticket.
    \end{itemize}
\end{enumerate}

\subsection{Decisiones de diseño reflejadas}
\begin{itemize}
  \item \textbf{Sin reintento automático post-ACK (LINK-03):} reenviar una venta acusada puede originar
    duplicación de cargos en la cuenta del tarjetahabiente, generando reclamos y sanciones comerciales.
  \item \textbf{Idempotencia a nivel de aplicación:} la reconciliación delegada en el host empareja la orden
    de venta con el registro oficial del hardware mediante identificadores unívocos.
\end{itemize}

\subsection{Guía de lectura}
Observe los cuadros destacados de advertencia en color rosa (\texttt{\#F2DEDE}) que señalan la zona de
incertidumbre transaccional, seguidos por el bloque alternativo (\texttt{alt/else}) final donde se
resuelve la consistencia de la venta.

\clearpage
\section{Diagrama 09: Despliegue de infraestructura y conectividad}

\begin{figure}[htbp]
  \centering
  \includegraphics[width=\textwidth]{despliegue.png}
  \caption{Diagrama de despliegue de infraestructura y conectividad.}
  \label{fig:despliegue}
\end{figure}

\subsection{Propósito}
El diagrama de despliegue (Figura~\ref{fig:despliegue}) muestra la topología de ejecución física y lógica
del SDK, detallando los controladores de hardware soportados a nivel de sistema operativo (Linux, macOS,
Windows), la emulación en red para entornos de desarrollo y el enlace con el switch adquirente.

\subsection{Topología y nodos de despliegue}
\begin{itemize}
  \item \textbf{Host / Computador de Caja:} equipo donde corre la aplicación de comercio (por ejemplo,
    \texttt{rework-sistema-facturacion} sobre Tauri) vinculada estáticamente con el crate \texttt{transbank-pos}.
    Interactúa con los controladores del kernel del sistema operativo:
    \begin{itemize}
      \item \textbf{Linux:} controlador CDC-ACM o FTDI (\texttt{/dev/ttyACM0}, \texttt{/dev/ttyUSB0})
        gobernado por reglas de \texttt{udev} que asignan permisos al grupo de usuarios \texttt{dialout}.
      \item \textbf{macOS:} subsistema IOKit gestionando dispositivos \texttt{/dev/cu.usbserial-*} o
        \texttt{/dev/cu.usbmodem*}.
      \item \textbf{Windows:} controlador nativo de módem serie USB (\texttt{usbser.sys}) o Virtual COM Port (VCP)
        asignando puertos COM correlativos (\texttt{COM3}, \texttt{COM4}).
    \end{itemize}
  \item \textbf{Entorno Local de Pruebas / CI:} nodo de desarrollo local o contenedor de integración continua
    donde se ejecuta el binario \texttt{tbk-pos-sim} escuchando en la interfaz de loopback \texttt{127.0.0.1:7777}.
  \item \textbf{Terminal Físico POS Transbank:} dispositivo hardware montado en caja (Verifone VX520,
    Ingenico Desk 3500) que integra un puente USB-UART comercial (FTDI, CP2102, CH340) operando a 115.200 baudios
    en POS Integrado y 19.200 baudios en Autoservicio.
  \item \textbf{Infraestructura Transbank:} switch adquirente centralizado conectado con el terminal a través
    de enlaces seguros Ethernet cableado, redes móviles 3G/4G o WiFi dedicado.
\end{itemize}

\subsection{Decisiones de diseño reflejadas}
\begin{itemize}
  \item \textbf{Restricción de red local en simulador (SEC-04):} el servidor del simulador enlaza por defecto
    únicamente en \texttt{127.0.0.1}, impidiendo que terminales no autorizados en la red corporativa simulen pagos falsos.
  \item \textbf{Independencia de plataforma en puerto serie (SAFE-03):} uso del crate \texttt{serialport}
    para garantizar compatibilidad transparente entre sistemas POS bajo Linux, Windows y macOS.
\end{itemize}

\subsection{Guía de lectura}
Las conexiones continuas representan interfaces físicas por cable serie/USB. Las líneas punteadas muestran
comunicaciones virtuales sobre TCP loopback. Las notas inferiores resumen los postulados de seguridad y aislamiento.

\clearpage
\section{Diagrama 10: Integración en el ecosistema raysac}

\begin{figure}[htbp]
  \centering
  \includegraphics[width=\textwidth]{ecosistema.png}
  \caption{Diagrama de integración en el ecosistema de retail y facturación raysac.}
  \label{fig:ecosistema}
\end{figure}

\subsection{Propósito}
El diagrama de integración (Figura~\ref{fig:ecosistema}) ilustra cómo \texttt{transbank-SDK} se inserta
de manera limpia y sin fricción dentro del portafolio más amplio de productos de software de retail,
pago y tributación desarrollados por el ecosistema raysac.

\subsection{Relaciones entre sistemas}
\begin{itemize}
  \item \textbf{Puntos de Venta (Retail):} el software de facturación \texttt{rework-sistema-facturacion}
    gestiona el carro de compras y los cobros de caja. Cuando el medio de pago es tarjeta, delega en
    \texttt{transbank-SDK} la orquestación del terminal.
  \item \textbf{Servicios Tributarios (\texttt{sii-sdk}):} una vez que \texttt{transbank-SDK} confirma
    la autorización bancaria mediante un \texttt{SaleResponse} aprobado, el software host invoca
    \texttt{sii-sdk} para generar y firmar digitalmente la Boleta Electrónica de Ventas y Servicios (DTE 39/41),
    remitiéndola a los servidores del SII.
  \item \textbf{Herramientas de Soporte y Calidad:} \texttt{tbk-pos-sim} permite certificar la aplicación de
    facturación en pipelines de CI desprovistos de hardware, mientras que \texttt{tbk-pos} (CLI) facilita
    la instalación inicial y auditoría en dependencias comerciales del cliente.
\end{itemize}

\subsection{Decisiones de diseño reflejadas}
\begin{itemize}
  \item \textbf{Biblioteca agnóstica del negocio (ADR-01, ADR-02):} \texttt{transbank-SDK} no contiene
    tipos de facturación, ni referencias al SII, ni esquemas de bases de datos de inventario. Es un
    componente de infraestructura reutilizable por cualquier software de punto de venta en Rust.
  \item \textbf{Desacoplamiento temporal y de red:} las transacciones financieras con Transbank y la emisión
    tributaria con el SII se ejecutan en fases lógicamente secuenciadas y atómicas.
\end{itemize}

\subsection{Guía de lectura}
Siga la numeración ordinal del 1 al 8 para rastrear el ciclo completo de una venta en tienda: desde la
intención de cobro, pasando por la validación bancaria en el POS, hasta el despacho final del documento
fiscal al Servicio de Impuestos Internos.

\clearpage
\section{Glosario de términos de modelado}

\begin{glosario}
  \termino{ACK}{Acknowledge (0x06). Carácter de control ASCII utilizado para acusar recibo conforme de una trama.}
  \termino{C4 Model}{Metodología de modelado de arquitectura de software en cuatro niveles de abstracción: Contexto, Contenedores, Componentes y Código.}
  \termino{ETX}{End of Text (0x03). Carácter de control ASCII que marca la finalización del payload de una trama.}
  \termino{LRC}{Longitudinal Redundancy Check. Checksum calculado mediante XOR acumulativo desde el primer byte del payload hasta ETX inclusive.}
  \termino{NAK}{Negative Acknowledge (0x15). Carácter de control ASCII emitido ante errores de verificación de checksum (LRC).}
  \termino{Newtype}{Patrón idiomático de Rust que envuelve un tipo primitivo en una tupla struct de un solo elemento para garantizar validaciones estrictas en tiempo de compilación.}
  \termino{POS}{Point of Sale. Terminal electrónico físico para captura de tarjetas bancarias.}
  \termino{Sans-IO}{Patrón de diseño arquitectónico que separa por completo los algoritmos de protocolo y parsing de cualquier operación de entrada/salida o reloj del sistema.}
  \termino{STX}{Start of Text (0x02). Carácter de control ASCII que delimita el inicio de una trama.}
  \termino{Tail Timeout}{Plazo de silencio (500 ms) empleado por la capa de enlace para detectar el fin de una respuesta compuesta por múltiples tramas sucesivas.}
\end{glosario}

\end{document}
